% 编辑转录副本；不是独立下载原件。来源：https://terrytao.wordpress.com/2016/05/18/a-symmetric-formulation-of-the-croot-lev-pach-ellenberg-gijswijt-capset-bound/

A capset in the vector space ${\bf F}_3^n$ over the finite field ${\bf F}_3$ of three elements is a subset $A$ of ${\bf F}_3^n$ that does not contain any lines $\{x,x+r,x+2r\}$, where $x,r\in{\bf F}_3^n$ and $r\ne0$.

The basic starting point is this: if $A$ is a capset, then one has the identity
\begin{equation}
\delta_{0^n}(x+y+z)=\sum_{a\in A}\delta_a(x)\delta_a(y)\delta_a(z)\tag{1}
\end{equation}
for all $(x,y,z)\in A^3$, where $\delta_a(x):=1_{a=x}$ is the Kronecker delta function, which we view as taking values in ${\bf F}_3$. Indeed, (1) reflects the fact that the equation $x+y+z=0$ has solutions precisely when $x,y,z$ are either all equal, or form a line, and the latter is ruled out precisely when $A$ is a capset.

To exploit (1), we will show that the left-hand side of (1) is ``low rank'' in some sense, while the right-hand side is ``high rank''.

Recall that a function $F:A\times A\to{\bf F}$ taking values in a field ${\bf F}$ is of rank one if it is non-zero and of the form $(x,y)\mapsto f(x)g(y)$ for some $f,g:A\to{\bf F}$, and that the rank of a general function $F:A\times A\to{\bf F}$ is the least number of rank one functions needed to express $F$ as a linear combination.

More generally, if $k\ge2$, we define the rank of a function $F:A^k\to{\bf F}$ to be the least number of ``rank one'' functions of the form
\[
(x_1,\ldots,x_k)\mapsto f(x_i)g(x_1,\ldots,x_{i-1},x_{i+1},\ldots,x_k)
\]
for some $i=1,\ldots,k$ and some functions $f:A\to{\bf F}$, $g:A^{k-1}\to{\bf F}$, that are needed to generate $F$ as a linear combination. For instance, when $k=3$, the rank one functions take the form $(x,y,z)\mapsto f(x)g(y,z)$, $(x,y,z)\mapsto f(y)g(x,z)$, $(x,y,z)\mapsto f(z)g(x,y)$, and linear combinations of $r$ such rank one functions will give a function of rank at most $r$.

It is a standard fact in linear algebra that the rank of a diagonal matrix is equal to the number of non-zero entries. This phenomenon extends to higher dimensions:
\begin{lemma}[1, Rank of diagonal hypermatrices]
Let $k\ge2$, let $A$ be a finite set, let ${\bf F}$ be a field, and for each $a\in A$, let $c_a\in{\bf F}$ be a coefficient. Then the rank of the function
\begin{equation}
(x_1,\ldots,x_k)\mapsto\sum_{a\in A}c_a\delta_a(x_1)\cdots\delta_a(x_k)\tag{2}
\end{equation}
is equal to the number of non-zero coefficients $c_a$.
\end{lemma}
\begin{proof}
We induct on $k$. As mentioned above, the case $k=2$ follows from standard linear algebra, so suppose now that $k>2$ and the claim has already been proven for $k-1$.

It is clear that the function (2) has rank at most equal to the number of non-zero $c_a$ (since the summands on the right-hand side are rank one functions), so it suffices to establish the lower bound. By deleting from $A$ those elements $a\in A$ with $c_a=0$ (which cannot increase the rank), we may assume without loss of generality that all the $c_a$ are non-zero. Now suppose for contradiction that (2) has rank at most $|A|-1$, then we obtain a representation
\begin{equation}
\begin{split}
&\sum_{a\in A}c_a\delta_a(x_1)\cdots\delta_a(x_k)\\
&\quad=\sum_{i=1}^k\sum_{\alpha\in I_i}f_{i,\alpha}(x_i)g_{i,\alpha}(x_1,\ldots,x_{i-1},x_{i+1},\ldots,x_k)
\end{split}\tag{3}
\end{equation}
for some sets $I_1,\ldots,I_k$ of cardinalities adding up to at most $|A|-1$, and some functions $f_{i,\alpha}:A\to{\bf F}$ and $g_{i,\alpha}:A^{k-1}\to{\bf R}$.

Consider the space of functions $h:A\to{\bf F}$ that are orthogonal to all the $f_{k,\alpha}$, $\alpha\in I_k$ in the sense that
\[
\sum_{x\in A}f_{k,\alpha}(x)h(x)=0
\]
for all $\alpha\in I_k$. This space is a vector space whose dimension $d$ is at least $|A|-|I_k|$. A basis of this space generates a $d\times|A|$ coordinate matrix of full rank, which implies that there is at least one non-singular $d\times d$ minor. This implies that there exists a function $h:A\to{\bf F}$ in this space which is nowhere vanishing on some subset $A'$ of $A$ of cardinality at least $|A|-|I_k|$.

If we multiply (3) by $h(x_k)$ and sum in $x_k$, we conclude that
\[
\begin{split}
&\sum_{a\in A}c_ah(a)\delta_a(x_1)\cdots\delta_a(x_{k-1})\\
&\quad=\sum_{i=1}^{k-1}\sum_{\alpha\in I_i}f_{i,\alpha}(x_i)\tilde g_{i,\alpha}(x_1,\ldots,x_{i-1},x_{i+1},\ldots,x_{k-1})
\end{split}
\]
where
\[
\begin{split}
&\tilde g_{i,\alpha}(x_1,\ldots,x_{i-1},x_{i+1},\ldots,x_{k-1})\\
&\quad:=\sum_{x_k\in A}g_{i,\alpha}(x_1,\ldots,x_{i-1},x_{i+1},\ldots,x_k)h(x_k).
\end{split}
\]
The right-hand side has rank at most $|A|-1-|I_k|$, since the summands are rank one functions. On the other hand, from induction hypothesis the left-hand side has rank at least $|A|-|I_k|$, giving the required contradiction.
\end{proof}

On the other hand, we have the following (symmetrised version of a) beautifully simple observation of Croot, Lev, and Pach:
\begin{lemma}[2]
On $({\bf F}_3^n)^3$, the rank of the function $(x,y,z)\mapsto\delta_{0^n}(x+y+z)$ is at most $3N$, where
\[
N:=\sum_{\substack{a,b,c\ge0:\ a+b+c=n\ b+2c\le2n/3}}\frac{n!}{a!b!c!}.
\]
\end{lemma}
\begin{proof}
Using the identity $\delta_0(x)=1-x^2$ for $x\in{\bf F}_3$, we have
\[
\delta_{0^n}(x+y+z)=\prod_{i=1}^n(1-(x_i+y_i+z_i)^2).
\]
The right-hand side is clearly a polynomial of degree $2n$ in $x,y,z$, which is then a linear combination of monomials
\[
x_1^{i_1}\cdots x_n^{i_n}y_1^{j_1}\cdots y_n^{j_n}z_1^{k_1}\cdots z_n^{k_n}
\]
with $i_1,\ldots,i_n,j_1,\ldots,j_n,k_1,\ldots,k_n\in\{0,1,2\}$ with
\[
i_1+\cdots+i_n+j_1+\cdots+j_n+k_1+\cdots+k_n\le2n.
\]
In particular, from the pigeonhole principle, at least one of $i_1+\cdots+i_n,j_1+\cdots+j_n,k_1+\cdots+k_n$ is at most $2n/3$.

Consider the contribution of the monomials for which $i_1+\cdots+i_n\le2n/3$. We can regroup this contribution as
\[
\sum_\alpha f_\alpha(x)g_\alpha(y,z)
\]
where $\alpha$ ranges over those $(i_1,\ldots,i_n)\in\{0,1,2\}^n$ with $i_1+\cdots+i_n\le2n/3$, $f_\alpha$ is the monomial
\[
f_\alpha(x_1,\ldots,x_n):=x_1^{i_1}\cdots x_n^{i_n}
\]
and $g_\alpha:{\bf F}_3^n\times{\bf F}_3^n\to{\bf F}_3$ is some explicitly computable function whose exact form will not be of relevance to our argument. The number of such $\alpha$ is equal to $N$, so this contribution has rank at most $N$. The remaining contributions arising from the cases $j_1+\cdots+j_n\le2n/3$ and $k_1+\cdots+k_n\le2n/3$ similarly have rank at most $N$ (grouping the monomials so that each monomial is only counted once), so the claim follows.
\end{proof}

Upon restricting from $({\bf F}_3^n)^3$ to $A^3$, the rank of $(x,y,z)\mapsto\delta_{0^n}(x+y+z)$ is still at most $3N$. The two lemmas then combine to give the Ellenberg--Gijswijt bound
\[
|A|\le3N.
\]
